%
% Exemplo LaTeX de monografia UNISINOS
%
% Elaborado com base nas orientações dadas no documento
% ``GUIA PARA ELABORAÇÃO DE TRABALHOS ACADÊMICOS''
% disponível no site da biblioteca da Unisinos.
% http://www.unisinos.br/biblioteca
%
% Os elementos textuais abaixo são apresentados na ordem em que devem
% aparecer no documento.  Repare que nem todos são obrigatórios - isso
% é devidamente indicado em cada caso.
%
% Comentários abaixo colocados entre aspas (`` '') foram
% extraídos diretamente do documento da biblioteca.
%
% Este documento é de domínio público.
%

%=======================================================================
% Declarações iniciais identificando a classe de documento e
% selecionando alguns pacotes adicionais.
%
% As opções disponíveis (separe-as com vírgulas, sem espaço) são:
% - twoside: Formata o documento para impressão frente-e-verso
%   (o default é somente-frente)
% - english,brazilian,french,german,etc.: idiomas usados no documento.
%   Deve ser colocado por último o idioma principal.
%=======================================================================
\documentclass[twoside,english,brazilian]{UNISINOSmonografia}
\usepackage[utf8]{inputenc} % charset do texto (utf8, latin1, etc.)
\usepackage[T1]{fontenc} % encoding da fonte (afeta a sep. de sílabas)
\usepackage{graphicx} % comandos para gráficos e inclusão de figuras
\usepackage{bibentry} % para inserir refs. bib. no meio do texto

%=======================================================================
% Escolha do sistema para geração de referências bibliográficas.
%
% O default é usar o estilo unisinos.bst.  Comente a definição abaixo
% e descomente a linha seguinte para usar o estilo do ABNTeX (é
% necessário ter esse pacote instalado).
%
% A vantagem do unisinos.bst é que ele permite o uso de um arquivo .bib
% seguindo as orientações tradicionais do BibTeX (veja essas orientações
% em http://ctan.tug.org/tex-archive/biblio/bibtex/contrib/doc/btxdoc.pdf).
% Entretanto, o estilo não suporta algumas citações mais exóticas como
% apud.  Para isso, use o ABNTeX, mas esteja ciente de que muitas de
% suas referências serão incompatíveis com os estilos tradicionais do
% BibTeX como plain, alpha, ieeetr, entre outros.
%=======================================================================
\usepackage[alf]{abntex2cite}
%\usepackage[alf]{abntcite}

%=======================================================================
% Dados gerais sobre o trabalho.
%=======================================================================
\autor{Hoffmann}{César}
\titulo{Não tem mais cardápio?}
\subtitulo{Entendendo o impacto da digitalização de bares e restaurantes entre 2020 e 2026 em Porto Alegre}
\orientador[PhD]{Lacerda}{Guilherme}
%\coorientador[Prof.~Dr.]{Lamport}{Leslie}

\unidade{Unidade Acadêmica de Graduação}
\curso{Curso de Bacharelado em Ciência da Computação}
\natureza{%
Monografia apresentada como requisito parcial para obtenção do título de Bacharel em Ciência da Computação, pelo Curso de Ciência da Computação da Universidade do Vale do Rio dos Sinos (UNISINOS)
}
\local{Porto Alegre}
\ano{2026}

% dados da ficha catalográfica
% (obrigatória somente para dissertações e teses)
%\cip{Dissertação (mestrado)}{004.732}
%\bibliotecario{Bibliotecária responsável: Fulana da Silva}{12/3456}

% cada palavra-chave deve ser fornecida duas vezes, uma em português e
% outra no idioma estrangeiro (na verdade, em tantos idiomas quantos se
% desejar).
\palavrachave{brazilian}{Cardápio digital}
\palavrachave{brazilian}{Transformação digital}
\palavrachave{brazilian}{Autoatendimento}
\palavrachave{brazilian}{QR Code}
\palavrachave{brazilian}{Bares e restaurantes}
\palavrachave{brazilian}{Estudo de caso}
\palavrachave{english}{Digital menu}
\palavrachave{english}{Digital transformation}
\palavrachave{english}{Self-ordering}
\palavrachave{english}{QR Code}
\palavrachave{english}{Bars and restaurants}
\palavrachave{english}{Case study}

%=======================================================================
% Início do documento.
%=======================================================================
\begin{document}
\capa
\folhaderosto
%\folhadeaprovacao % não deve ser incluída nos TCCs

%=======================================================================
% Dedicatória (opcional).
%
% O texto é normalmente colocado na parte de baixo da página, alinhado
% à direita.  Mas a formatação é basicamente livre.  Só não se escreve
% a palavra 'dedicatória'.
%=======================================================================
\begin{dedicatoria}
A todos os pilares que sustentaram esta ponte e tornaram a jornada possível.\\[1ex] % quebra a linha dando um espaçamento maior
Olavo, Regina, Cristina, Marina, Vanessa e Fernanda\\[4ex] % quebra a linha dando um espaçamento maior

\begin{itshape} % faz o texto ficar em itálico
Eu sou o conjunto dessas pessoas\\
- e felizmente sou gordo o bastante pra que todas caibam no meu corpinho.\\
% \vspace{0.2cm}
\end{itshape}
--- \textsc{José Eugênio Soares} % \textsc é o "small caps"
\end{dedicatoria}

%=======================================================================
% Agradecimentos (opcional).
%=======================================================================
% \begin{agradecimentos}
% Obrigado!
% \end{agradecimentos}

%=======================================================================
% Epígrafe (opcional).
%
% ``[...] o autor apresenta uma citação, seguida de indicação de autoria,
% relacionada com a matéria tratada no corpo do trabalho. Podem, também,
% constar epígrafes nas folhas de aberturas das seções primárias.''
%=======================================================================
\begin{epigrafe}
``\textit{Ninguém abre um livro sem que aprenda alguma coisa}''.\\
(Anônimo)
\end{epigrafe}

%=======================================================================
% Resumo em Português.
%
% A recomendação é para 150 a 500 palavras.
%=======================================================================
\begin{abstract}
Este projeto de pesquisa investiga como bares e restaurantes de Porto Alegre que adotaram cardápios digitais com autoatendimento desde 2020 percebem os benefícios, riscos e desafios técnicos e operacionais dessa transição, bem como a forma com que seus clientes vivenciam essa nova interação. A pesquisa adota uma abordagem qualitativa-quantitativa por meio de um estudo de caso múltiplo envolvendo de dois a três estabelecimentos com perfis distintos de adoção, combinado com um questionário aplicado a frequentadores do setor. Espera-se contribuir para a compreensão da transformação digital em pequenos e médios estabelecimentos do food service e fornecer subsídios para o projeto de soluções de cardápio digital mais alinhadas às necessidades de gestores e clientes. Este documento corresponde à entrega do TCC I e apresenta a delimitação do tema, a fundamentação teórica preliminar, a metodologia proposta e o cronograma de execução da etapa de coleta e análise prevista para o TCC II.
\end{abstract}

%=======================================================================
% Resumo em língua estrangeira (obrigatório somente para teses e
% dissertações).
%
% O idioma usado aqui deve necessariamente aparecer nos parâmetros do
% \documentclass, no início do documento.
%=======================================================================
\begin{otherlanguage}{english}
\begin{abstract}
This research project investigates how bars and restaurants in Porto Alegre that adopted digital self-ordering menus since 2020 perceive the benefits, risks and technical or operational challenges of this transition, as well as how their customers experience this new form of interaction. The research adopts a qualitative-quantitative approach through a multiple case study involving two to three establishments with different adoption profiles, combined with a survey applied to customers. The expected contribution is a better understanding of digital transformation in small and medium-sized food-service establishments and inputs for the design of digital menu solutions that are better aligned with the needs of both managers and customers. This document corresponds to the first part of the final undergraduate work (TCC I) and presents the scope, preliminary theoretical foundation, proposed methodology and execution timeline for the data collection and analysis stage planned for TCC II.
\end{abstract}
\end{otherlanguage}

%=======================================================================
% Lista de Figuras (opcional).
%=======================================================================
\listoffigures

%=======================================================================
% Lista de Tabelas (opcional).
%=======================================================================
\listoftables

%=======================================================================
% Lista de Abreviaturas (opcional).
%
% Deve ser passada como parâmetro a maior das abreviaturas utilizadas.
%=======================================================================
% Lista de abreviaturas desabilitada — não há abreviaturas relevantes neste TCC.
% \begin{listadeabreviaturas}{seg., segs.}
% \item[ampl.] ampliado, -a
% \end{listadeabreviaturas}

%=======================================================================
% Lista de Siglas (opcional).
%
% Deve ser passada como parâmetro a maior das siglas utilizadas.
%=======================================================================
\begin{listadesiglas}{ABRASEL}
\item[ABNT] Associação Brasileira de Normas Técnicas
\item[ABRASEL] Associação Brasileira de Bares e Restaurantes
\item[CEP] Comitê de Ética em Pesquisa
\item[KDS] Kitchen Display System
\item[LGPD] Lei Geral de Proteção de Dados Pessoais
\item[POA] Porto Alegre
\item[QR] Quick Response (code)
\item[SEBRAE] Serviço Brasileiro de Apoio às Micro e Pequenas Empresas
\item[TAM] Technology Acceptance Model
\item[TCC] Trabalho de Conclusão de Curso
\item[TCLE] Termo de Consentimento Livre e Esclarecido
\item[UNISINOS] Universidade do Vale do Rio dos Sinos
\item[UTAUT] Unified Theory of Acceptance and Use of Technology
\item[UX] User Experience
\end{listadesiglas}

%=======================================================================
% Lista de Símbolos (opcional).
%
% Deve ser passado o maior (mais largo) dos símbolos utilizados.
%=======================================================================
% Lista de símbolos desabilitada — não há símbolos relevantes neste TCC.
% \begin{listadesimbolos}{Ca}
% \item[Ca] Cálcio
% \end{listadesimbolos}

%=======================================================================
% Sumário
%=======================================================================
\tableofcontents

%=======================================================================
% 1 Introdução
%=======================================================================
\chapter{Introdução}

% TODO: parágrafo introdutório de abertura do capítulo (1--2 parágrafos),
% conectando o leitor ao assunto antes de entrar nas subseções.

\section{Tema}
\label{sec:tema}

% TODO: apresentar o tema em ~1 página.
% - Contexto da transformação digital no setor de bares e restaurantes.
% - Surgimento e expansão dos cardápios digitais (QR Code, web apps) após 2020.
% - Por que esse fenômeno chama atenção em Porto Alegre.

\section{Delimitação do tema}
\label{sec:delimitacao}

% TODO: delimitar escopo em quatro eixos:
% - Geográfico: Porto Alegre/RS, 2--3 bairros (Cidade Baixa, Moinhos, Centro).
% - Temporal: adoções entre mar/2020 e 2026, foco antes/depois da pandemia.
% - Funcional: cardápio digital + autoatendimento via QR/web
%   (NÃO inclui KDS, delivery, integração com iFood, ERP).
% - Sujeitos: 1--3 estabelecimentos como casos + frequentadores via survey.

\section{Problema}
\label{sec:problema}

% TODO: formular como pergunta única:
% Como bares e restaurantes de Porto Alegre que adotaram cardápios digitais
% com autoatendimento desde 2020 percebem os benefícios, riscos e desafios
% dessa transição, e como seus clientes vivenciam essa interação?

\section{Objetivos}
\label{sec:objetivos}

\subsection{Objetivo geral}
\label{sec:obj-geral}

% TODO: analisar, por meio de estudo de caso múltiplo, como o uso de
% cardápios digitais com autoatendimento impacta a operação de bares
% e restaurantes de Porto Alegre e a experiência de seus clientes.

\subsection{Objetivos específicos}
\label{sec:obj-especificos}

\begin{alineas}
	\item caracterizar o estado da arte sobre cardápios digitais, QR menus e self-ordering em food service;
	\item levantar, com gestores de 1 a 3 estabelecimentos, as motivações, custos e desafios técnicos da adoção;
	\item coletar, via questionário, a percepção de frequentadores quanto a usabilidade, confiança e preferências;
	\item confrontar a percepção dos estabelecimentos com a percepção dos clientes, destacando convergências e tensões;
	\item discutir aspectos técnicos transversais: UX em mobile web, LGPD, acessibilidade e dependência de conectividade.
\end{alineas}

\section{Justificativa}
\label{sec:justificativa}

% TODO: justificar em dois eixos:
% - Relevância prática: setor altamente impactado pela pandemia; decisões
%   de digitalização envolvem investimentos relevantes para PMEs.
% - Relevância para a Ciência da Computação: o objeto é tecnológico
%   (web mobile, QR Code, UX, LGPD); os achados podem orientar requisitos
%   de software, padrões de UX e práticas de privacidade para o setor.

%=======================================================================
% 2 Fundamentação Teórica
%=======================================================================
\chapter{Fundamentação Teórica}

% TODO: parágrafo de abertura indicando como as seções a seguir se
% articulam e qual o caminho conceitual percorrido.

\section{Transformação digital no setor de food service}
\label{sec:td-foodservice}

% TODO: conceito de transformação digital; impacto específico no setor
% de bares e restaurantes; aceleração causada pela pandemia de COVID-19.
% Apoiar com autores de TD em PMEs e relatórios setoriais (ABRASEL).

\section{Cardápios digitais e QR Code}
\label{sec:cardapios-qr}

% TODO: definição e tipologia (PDF estático, web app, app nativo);
% histórico breve da tecnologia QR Code aplicada ao varejo; arquiteturas
% comuns (link -> web app -> backend de cardápio).

\section{Autoatendimento e experiência do usuário em mobile}
\label{sec:autoatendimento-ux}

% TODO: conceito de self-ordering; heurísticas de usabilidade aplicadas
% a cardápios mobile (Nielsen); acessibilidade (WCAG); fricções típicas
% (QR não escaneia, conexão fraca, navegação confusa).

\section{LGPD aplicada a restaurantes}
\label{sec:lgpd}

% TODO: panorama da LGPD; dados pessoais coletados via cardápio digital
% (e-mail, telefone, geolocalização, comportamento de navegação); base
% legal aplicável; obrigações do controlador.

\section{Trabalhos relacionados}
\label{sec:trabalhos-relacionados}

% TODO: 4--8 trabalhos (TCCs, dissertações, artigos) que tratem de
% digitalização em restaurantes, QR menus ou adoção de tecnologia em PMEs.
% Encerrar com quadro comparativo (foco, método, recorte, lacuna).

%=======================================================================
% 3 Metodologia
%=======================================================================
\chapter{Metodologia}

% TODO: parágrafo de abertura indicando que a pesquisa é qualitativa-quantitativa
% e que este capítulo descreve a abordagem, os instrumentos e os procedimentos.

\section{Caracterização da pesquisa}
\label{sec:caracterizacao}

% TODO: classificar quanto à natureza (aplicada), abordagem
% (qualitativa-quantitativa / métodos mistos), objetivos
% (descritiva-exploratória) e procedimentos (estudo de caso múltiplo + survey).
% Apoiar em Yin (estudo de caso) e Gil/Creswell (métodos).

\section{Etapas da pesquisa}
\label{sec:etapas}

% TODO: descrever fluxo da fundamentação até a defesa.
% Sugestão: incluir uma figura com fluxograma das etapas TCC I -> TCC II.

\section{Instrumentos de coleta}
\label{sec:instrumentos}

\subsection{Entrevista semiestruturada com gestores}
\label{sec:entrevista}

% TODO: descrever a entrevista (45--60 min, gravada mediante TCLE),
% a estrutura modular do roteiro (Blocos 0--5) e a relação com os
% perfis A/B/C dos estabelecimentos. Remeter ao Apêndice A.

\subsection{Questionário com frequentadores}
\label{sec:questionario}

% TODO: descrever o survey (Google Forms, 15--20 perguntas, meta de 80+
% respostas), canais de distribuição (redes sociais, grupos, opcionalmente
% QR Code nos estabelecimentos) e estrutura em 4 seções. Remeter ao Apêndice B.

\section{Critérios de seleção dos estabelecimentos}
\label{sec:criterios}

% TODO: amostragem intencional com perfis
% A (pré-2020 que adotou), B (pré-2020 que resistiu) e C (pós-2020 nascido digital).
% Critérios de inclusão (porte pequeno/médio, tomador de decisão acessível)
% e de exclusão (redes franqueadas, fine dining, dark kitchens).

\section{Plano de análise}
\label{sec:analise}

% TODO: análise de conteúdo (Bardin) para o material qualitativo das
% entrevistas; estatística descritiva (frequências, médias, cruzamentos)
% para o survey. Plano de triangulação entre as duas fontes.

\section{Aspectos éticos}
\label{sec:etica}

% TODO: TCLE com participantes (Apêndice C); aderência à LGPD;
% Termo de Confidencialidade (modelo UNISINOS) com gestores;
% Termo de Autorização de Uso de Imagem quando aplicável;
% discussão sobre necessidade ou não de submissão ao CEP.

%=======================================================================
% 4 Cronograma
%=======================================================================
\chapter{Cronograma}

% TODO: parágrafo introdutório descrevendo o planejamento temporal
% da execução do TCC II (2026/2) e os marcos principais.

\begin{table}
	\caption{Cronograma de execução do TCC II em 2026/2}
	\label{tab:cronograma}
	\centering%
	\begin{minipage}{.98\textwidth}
		\begin{tabular*}{\textwidth}{l@{\extracolsep{\fill}}cccccc}
			\hline
			\textbf{Atividade} & \textbf{Ago} & \textbf{Set} & \textbf{Out} & \textbf{Nov} & \textbf{Dez} & \textbf{Defesa}\\
			\hline
			Piloto dos instrumentos          & X &   &   &   &   &  \\
			Recrutamento de estabelecimentos & X & X &   &   &   &  \\
			Entrevistas com gestores         &   & X & X &   &   &  \\
			Aplicação do questionário        &   & X & X & X &   &  \\
			Transcrição e codificação        &   &   & X & X &   &  \\
			Análise dos dados                &   &   &   & X & X &  \\
			Redação do artigo estendido      &   &   & X & X & X &  \\
			Revisão final                    &   &   &   &   & X &  \\
			Defesa                           &   &   &   &   &   & X\\
			\hline
		\end{tabular*}
		\fonte{Elaborado pelo autor.}
	\end{minipage}
\end{table}

%=======================================================================
% Referências
%=======================================================================
\bibliography{exemplo}

%=======================================================================
% Apêndices
%=======================================================================
\appendix

%-----------------------------------------------------------------------
% Apêndice A --- Roteiro de entrevista com gestores
%-----------------------------------------------------------------------
\chapter{Roteiro de entrevista com gestores}
\label{ap:roteiro}

% TODO: roteiro modular completo da entrevista semiestruturada
% (Blocos 0 a 5), conforme attachments/plano-tcc1.md (seção 6.1).
% Versão definitiva será produzida na fase de instrumentos.

%-----------------------------------------------------------------------
% Apêndice B --- Questionário com frequentadores
%-----------------------------------------------------------------------
\chapter{Questionário com frequentadores}
\label{ap:questionario}

% TODO: questionário completo (15--20 perguntas), aplicado via Google Forms,
% organizado em quatro seções (Perfil; Experiência com cardápios digitais;
% Aspectos técnicos; Visão geral). Ver attachments/plano-tcc1.md (seção 6.2).

%-----------------------------------------------------------------------
% Apêndice C --- Termos de consentimento e autorização
%-----------------------------------------------------------------------
\chapter{Termos de consentimento e autorização}
\label{ap:termos}

% TODO: incluir nesta ordem:
% - TCLE (Termo de Consentimento Livre e Esclarecido) --- a ser redigido.
% - Termo de Confidencialidade (modelo UNISINOS).
% - Termo de Autorização de Uso de Imagem (modelo UNISINOS), se aplicável.
% - Autorização de Uso de Obra Fotográfica (modelo UNISINOS), se aplicável.
% Modelos em attachments/autorizações/.

\end{document}
